\documentclass{standalone}
\usepackage{circuitikz}

\begin{document}

\begin{circuitikz}

% Línea principal con resistencias
\draw
(0,0) node[left]{$T_j$}
to[short, o-] (1,0)
to[R] (3,0)
to[R] (5,0)
to[R] (7,0)
to[R] (9,0)
to[short] (9,0) node[right]{$T_{case}$};

% Capacitores del modelo
\draw (1, 0) to[C] (1,-2);
\draw (3, 0) to[C] (3,-2);
\draw (5, 0) to[C] (5,-2);
\draw (7, 0) to[C] (7,-2);


% Conexión a T_case
\draw (9, 0) to[short] (9, -2)
node[ground]{}
node[below]{};

\draw (1, -2) to[short] (1, -2)
node[ground]{}
node[below]{};

\draw (3, -2) to[short] (3, -2)
node[ground]{}
node[below]{};

\draw (5, -2) to[short] (5, -2)
node[ground]{}
node[below]{};

\draw (7, -2) to[short] (7, -2)
node[ground]{}
node[below]{};



\end{circuitikz}

\end{document}